% Main Preamble

\documentclass[12pt, letterpaper]{article}
\title{Nucleofile Manual\\[0.12in] \large Version 0.1.1}
\author{Caleb Griffin}
\date{June, 2026}

\usepackage{
	xcolor, 
	graphicx, 
	titlesec, 
	geometry, 
	listings, 
	amsmath,
	hyperref
}
\geometry{bottom=1.0in}

\pagecolor{white}

\hypersetup{
	colorlinks=true,
	linkcolor=black,
	filecolor=magenta,
	urlcolor=blue
}

\begin{document}
\maketitle
\thispagestyle{empty}

\pagebreak
\tableofcontents
\pagebreak

\section{Introduction}
Nucleofile is meant to serve as a general purpose software aid to chemists
and all who need scientific programs. Nucleofile will always be free and
and open source.

\section{User Guide}
Nucleofile is designed to be flexible and work with you through many
interfaces. The three primary ways to interact with nucleofile are
with the GUI, TUI, and the interpreter. Project files are saved as ``.ncfl''
extension. The app can also be controled through a script languaged called ``aangstrom'' which uses the ``.aang'' suffix for files.

\subsection{Flags}
\begin{itemize}
\item -f/--file Pass a file or diretory to nucleofile.
\item -d/--debug
\item -g/--gui Open nucleofile without opening the GUI.
\end{itemize}

\subsection{TUI Guide}
To use the TUI, open nucleofile with ``ncfl -t''.

\section{Math}
There are multiple constants used in nucleofile, these are stated here
 with bold script as textbf{constant}.

\subsection{Determining the area under a function}
If you want to determine \(\int_a^b x^2\,dx\) for your system.

\end{document}
